\section{Laid-Line Detection and Evaluation}
\label{sec:detection}

% PAGE TARGET: ~7 pages.
% Day 4 (5/28). Narrative: Initial Gabor / multi-phi attempt -> revised
% spectral-power dual-band -> polarity-flip discovery -> evaluation
% metrics. Andrew-style: lots of figures, few displayed equations,
% honest failure narrative.

\subsection{Initial Approach: Gabor Filter Bank}
\label{sec:detection:gabor}
% NARRATIVE: Why Gabor was tried first; what worked; what failed.
%
% Cover:
% - Gabor is the natural starting point because the enhanced greyscale
%   of Section 2 has a clear dominant orientation and a near-sinusoidal
%   profile along the line normal.
% - A Gabor bank covering periods {6, 8, ..., 40} px at the patch level
%   was applied, and the patchwise dominant period was aggregated into
%   a folio-level distribution.
% - On synthetic phantoms with clean signals and orientations close to
%   the bank's central wavelength the method recovered the true period
%   to within a few per cent.
% - On real folios two failure modes emerged: (i) the dominant period
%   landed close to the edge of the bank for high- or low-density
%   folios, producing an octave aliasing of the kind documented in
%   Section 4.X; (ii) the patch-level estimate was unstable on folios
%   with heavy ink, producing a long tail in the aggregated
%   distribution that biased the folio-level median.
% TODO: Figure 3.1 -- Gabor filter bank diagram (one row of filters
% at the bank wavelengths), and an example aggregated distribution
% on a phantom and on a real folio.

\subsection{Multi-Phi Detector}
\label{sec:detection:multiphi}
% NARRATIVE: A different angle on the same data; why we kept it as a
% cross-check rather than as the primary detector.
%
% Cover:
% - The multi-phi detector reuses the patch direction field of
%   Section 2.3, slicing the page into bands perpendicular to the
%   local laid direction and running a one-dimensional peak-find
%   along each band.
% - Pros: respects local orientation variation; no fixed period bank.
% - Cons: peak-find is sensitive to ink-induced spurious peaks and to
%   the choice of smoothing kernel; produces a per-band estimate
%   rather than a single folio-level number.
% - We retained this detector for cross-validation and for the
%   per-segment FWHM estimate of Section 3.5 (wire width), but did not
%   adopt it as the primary period estimator.

\subsection{Revised Detector: Spectral Power with Dual-Band Cross-Validation}
\label{sec:detection:spectral}
% NARRATIVE: Why we replaced projection autocorrelation with spectral
% power; what dual-band cross-validation buys us. This is the central
% design move of the project and should get the most narrative weight.
%
% Cover:
% - Earlier versions of the auto-direction detector used the
%   autocorrelation of the column-projection signal to score
%   candidate orientations. On phantoms this worked, but on noisy
%   real folios the autocorrelation peak at the true period was
%   often only marginally larger than a peak at the noise correlation
%   length, producing high-confidence wrong answers.
% - The replacement (commit caecc39) substitutes the spectral power
%   in a narrow window around the candidate frequency. Spectral
%   power gives a flatter null distribution under noise, so the
%   contrast between the true peak and the noise floor improves.
% - Dual-band cross-validation runs the same scoring on two
%   non-overlapping horizontal bands of the page. An orientation is
%   accepted only if both bands report a peak at the same period;
%   single-band peaks (typically caused by an ink patch or a fold)
%   are rejected. This eliminates the false positives produced by
%   the older version on three of the nine benchmark folios.
% TODO: Figure 3.2 -- side-by-side scoring landscapes for
% autocorrelation vs spectral power on a real folio; show the
% improved peak contrast.
% TODO: Figure 3.3 -- dual-band cross-validation cartoon: same folio,
% two horizontal bands, both spectra shown, peak alignment circled.

\subsection{Polarity Flip for Transmitted-Light Samples}
\label{sec:detection:polarity}
% NARRATIVE: A small but real bug found by working on a TX folio.
%
% Cover:
% - The peak-finding code assumed that laid wires appear as local
%   maxima of the enhanced greyscale, which is true for reflective
%   acquisition (wires catch grazing light and look brighter than
%   the troughs).
% - For transmitted-light folios this is inverted: wires absorb more
%   light and look darker than the troughs. The very first TX folio
%   was therefore reported with a period twice the true value.
% - Fix (commit 38c37f3): a per-folio polarity flag
%   wire\_is\_darker that toggles a sign in the peak-finder. The
%   flag is set in the per-folio configuration file rather than
%   detected automatically; an automatic detector is discussed
%   in the future-work section.

\subsection{Evaluation Metrics}
\label{sec:detection:evaluation}
% Five metrics, each a JSON file plus a companion PNG. Describe each in
% one short paragraph, give the formula inline where possible.
%
% Cover:
% (a) Harmonic fit quality (fit_quality.json). $R^2$ of the
%     least-squares fit of the projection signal to a sum of
%     harmonics up to $k = 4$ of the detected period. A low $R^2$
%     indicates that the detected period does not in fact explain
%     the bulk of the periodic content.
%
% (b) Interval distribution (interval_distribution.json). The
%     empirical distribution of peak-to-peak intervals over the
%     detected lines; reported as mean, median and a count of the
%     peaks $n_\text{peaks}$. Useful for spotting bimodal
%     distributions caused by uneven wires.
%
% (c) Self-contrast z-score (self_contrast.json). The scoring
%     function used by the detector is recomputed on a randomly
%     rotated copy of the same image; the rotation breaks any true
%     periodicity, so the rotated score is a null sample.
%     $z = (s_{\mathrm{true}} - \mu_{\mathrm{null}}) /
%          \sigma_{\mathrm{null}}$ over 50 rotations gives a per-folio
%     significance level.
%
% (d) Split-half stability (split_half_stability.json). The folio is
%     split into a top half and a bottom half; the detector is run on
%     each half; the difference between the two estimates is reported
%     as a stability figure. Large split-half differences flag folios
%     with strong density variation across the page.
%
% (e) Wire width statistics (wire_width_stats.json). For each
%     detected line segment, the full-width-at-half-maximum of the
%     fitted profile is extracted; the per-folio median and a
%     bootstrap 95\% confidence interval are reported.
%
% Together these five metrics form the per-folio quality vector that
% the codicologist consumes alongside the headline density figure;
% Section 4.X shows the five metrics for each benchmark folio.
